\documentclass[a4paper]{article}

\usepackage[top=2.54cm,bottom=2.54cm,left=3.18cm,right=3.18cm]{geometry}
\usepackage{fontspec}
\usepackage{xeCJK}
\usepackage{setspace}
\usepackage{indentfirst}
\usepackage{enumitem}
\usepackage{hyperref}

\newcommand{\SongtiName}{Songti SC}
\IfFontExistsTF{Songti SC}{%
  \renewcommand{\SongtiName}{Songti SC}%
}{%
  \IfFontExistsTF{STSong}{%
    \renewcommand{\SongtiName}{STSong}%
  }{%
    \IfFontExistsTF{SimSun}{%
      \renewcommand{\SongtiName}{SimSun}%
    }{%
      \renewcommand{\SongtiName}{Noto Serif CJK SC}%
    }%
  }%
}

\setmainfont{Times New Roman}
\setCJKmainfont[AutoFakeBold=2.2]{\SongtiName}
\XeTeXlinebreaklocale "zh"
\XeTeXlinebreakskip = 0pt plus 1pt
\hypersetup{hidelinks}

\pagestyle{plain}
\setlength{\parindent}{2em}
\setlength{\parskip}{0pt}
\setstretch{1.0}

\newcommand{\exptitle}[1]{%
  \vspace*{7.8pt}
  \begin{center}
    {\fontsize{18pt}{21.6pt}\selectfont #1\par}
  \end{center}
  \vspace{7.8pt}
}

\newcommand{\studentinfo}{%
  \noindent 姓名：\underline{\makebox[2.65cm][c]{}}\quad
  学号：\underline{\makebox[2.65cm][c]{}}\par
  \vspace{7.8pt}
}

\newcommand{\reportsection}[1]{%
  \vspace{7.8pt}
  {\noindent\fontsize{14pt}{16.8pt}\selectfont #1\par}
  \vspace{7.8pt}
}

\newcommand{\reportsectionbold}[1]{%
  \vspace{7.8pt}
  {\noindent\fontsize{14pt}{16.8pt}\selectfont\bfseries #1\par}
  \vspace{7.8pt}
}

\newcommand{\normalpara}[1]{%
  \noindent #1\par
}

\newcommand{\indentpara}[1]{%
  \noindent\hspace*{2em}#1\par
}

\newlist{reportitems}{enumerate}{1}
\setlist[reportitems]{%
  label=\arabic*),
  leftmargin=42pt,
  labelwidth=21pt,
  labelsep=0pt,
  itemindent=0pt,
  align=left,
  topsep=0pt,
  partopsep=0pt,
  parsep=0pt,
  itemsep=0pt
}

\begin{document}
\fontsize{10.5pt}{15.6pt}\selectfont

\exptitle{实验一 XXX}

\studentinfo

\reportsectionbold{【问题描述】}
\indentpara{陈述本实验的任务，强调要解决的问题是什么？明确规定：输入的形式和输入值的范围；输出的形式；系统所能达到的功能；测试数据。}

\reportsection{【基本思路】}
\indentpara{写出存储结构，主要算法的基本思想（自然语言或类高级语言描述）。}
\indentpara{每个操作及模块的伪码算法。列出每个过程或函数所调用关系，可以通过层次图表达。}

\reportsection{【实验结果及分析】}
\normalpara{1. 实验结果： 请描述实验过程，包括编程环境、测试数据等，用图、表等方式，呈现实验结果。}

\indentpara{1.1 实验过程描述}

\begin{reportitems}
  \item 本实验用 XX 语言开发，编译环境为 XX，运行机器为本人笔记本电脑或 XXXX。由 X 人 设计方案，本人独立完成编码和测试。
  \item 测试数据选取思路为：典型，XXX，XXX。累计选取 X 组数据进行测试。
  \item XXX 部分的逻辑是测试重点。因为 XXX
\end{reportitems}

\noindent\hspace*{2em}1.2测试结果 用图、表等方式，呈现实验结果。\par

\normalpara{2. 结果及方案分析}

\indentpara{2.1本实验的重点及难点}
\indentpara{本实验设计的重点为 XXXX。 本实验设计的难点为 XXXX。}

\indentpara{2.2本实验设计的优点及可以改进的方向}
\indentpara{实验完全或部分实现了预计功能。}
\indentpara{设计得较为充分之处为 XX}
\indentpara{设计的可优化处包括：XXXX}

\normalpara{3. 还有哪些实现方法？谈谈你的个人观点。（思政目标）}
\indentpara{}

\normalpara{4. 附件：源程序代码}
\indentpara{}

\end{document}
